% Generated draft: requires mathematical and source review.
\documentclass[11pt]{article}
\usepackage[margin=25mm]{geometry}
\usepackage{fontspec}
\setmainfont{Latin Modern Roman}
\usepackage{amsmath,amssymb,mathtools,unicode-math}
\setmathfont{Latin Modern Math}
\usepackage{enumitem,needspace,xurl,hyperref}
\setlength{\parindent}{0pt}
\setlength{\parskip}{5pt}
\setlength{\emergencystretch}{3em}
\newcommand{\N}{\mathbb{N}}
\newcommand{\Z}{\mathbb{Z}}
\newcommand{\Q}{\mathbb{Q}}
\newcommand{\R}{\mathbb{R}}
\newcommand{\C}{\mathbb{C}}
\newcommand{\F}{\mathbb{F}}
\newcommand{\A}{\mathbb{A}}
\newcommand{\PP}{\mathbb{P}}
\newcommand{\E}{\mathbb{E}}
\newcommand{\eps}{\varepsilon}
\newcommand{\dd}{\,\mathrm{d}}
\begin{document}
\section*{UnsolvedMath ID: 800468}
\textbf{Source problem code:} catalogue entry
\subsection{Definitions and mathematical statement}
Consider a Newton equation with one configuration variable and periodic time dependence,
\[
\ddot u=-\partial_u V(u,t),\qquad L(u,\dot u,t)=\frac12\dot u^2-V(u,t).
\]
Let \(V:\R^2\to\R\) be periodic in both variables: for some periods \(P_u,P_t>0\),
\[
V(u+P_u,t)=V(u,t),\qquad V(u,t+P_t)=V(u,t).
\]
The source does not specify these periods or the regularity class of \(V\); sufficient regularity for the stated differential equation and flow is understood but requires clarification. With momentum \(p=\dot u\), the Hamiltonian is
\[
H(u,p,t)=\frac12p^2+V(u,t),
\]
and the extended equations are
\[
\dot u=p,\qquad \dot p=-\partial_uV(u,t),\qquad \dot t=1.
\]
Here \(u\) and \(t\) are taken modulo their respective periods, giving the extended phase space identified in the source as \(T^*\mathbb S^1\times\mathbb S^1\).

The source asserts that this flow is totally integrable on the entire phase space if and only if \(V\) is independent of \(u\). It also asserts that potentials of the form
\[
V(u,t)=W(nu+mt),\qquad n\in\Z,\quad m\ne0,
\]
admit an integral of motion quadratic in \(p\), and that their systems are consequently totally integrable near infinity, expressed there as \(|p|>c\) for sufficiently large \(c\). The notation \(W\) distinguishes the one-variable profile from the two-variable potential. The profile and parameters must satisfy the original double-periodicity requirement; the source does not specify the domain of \(m\), the period of \(W\), or the precise compatibility conditions.

\textbf{Question (Misha Bialy, Problem 3).} Are there other doubly periodic potentials \(V\), beyond the displayed family, for which the associated system is totally integrable at infinity, that is, for all sufficiently large \(|p|\)?

\textbf{Unresolved formulation qualifications.} The source does not define ``totally integrable,'' either globally or at infinity. In particular, it does not state the required invariant structures or their regularity, or whether the neighborhood at infinity must itself be invariant. Thus the question above retains the source terminology and is not a uniquely specified mathematical predicate until these choices are supplied. The equivalence relation, if any, under which potentials count as ``other'' is also unspecified.
\subsection{Short English statement}
For a Newton equation with a potential periodic in both position and time, the source asks whether total integrability for sufficiently large momentum occurs for potentials beyond the exhibited family \(V(u,t)=W(nu+mt)\), with \(n\in\Z\) and \(m\ne0\). The precise meaning of total integrability and several regularity and periodicity conventions remain unspecified. The supplied registry labels this problem open; that label is not an independently verified mathematical assessment.
\subsection{Sources}
\begin{itemize}
\item Local UnsolvedMath registry, permanent ID \(800468\), exact source code \texttt{catalogue entry}, title: ``Billiards and Quantitative Symplectic Geometry --- Misha Bialy --- Problem 3.'' The supplied record labels the problem \texttt{open} and records proposed year \(2026\).
\item Albach, Barraud et al., \emph{Open Problems in Billiards and Quantitative Symplectic Geometry} (2026), source item ``Misha Bialy --- Problem 3,'' \url{https://arxiv.org/abs/2602.12896}. The registry attributes the collection to Albach, Barraud et al.; the supplied background identifies Misha Bialy as the contributor.
\item Supplied provenance records access on \texttt{2026-10-16} and extraction method \texttt{arxiv\_latex}. This access information is supplied provenance and was not independently verified for this draft. The complete supplied record is the evidence used here; the linked source has not been visited for this draft.
\end{itemize}
\end{document}
